% ─────────────────────────────────────────────────────────────────────────────
\section{ELS~--- Les fl\`eches (\S{}7.2)}
% ─────────────────────────────────────────────────────────────────────────────

\begin{maitre}
Apr\`es avoir v\'erifi\'e que la poutre ne casse pas (ELU),
que reste-t-il \`a v\'erifier~?
\end{maitre}

\begin{apprenti}
Qu'elle ne fl\'echit pas trop~? Une poutre qui tient mais qui fait
une grande fl\`eche visible, c'est probl\'ematique pour les rev\^etements
et d\'esagr\'eable visuellement.
\end{apprenti}

\begin{maitre}
Exactement. C'est l'objet des \'Etats Limites de Service (ELS).
Ici on ne v\'erifie plus la r\'esistance, mais la \emph{d\'eformation}.
\end{maitre}

\subsection{Fl\`eche instantan\'ee}

Pour une poutre simplement appuy\'ee sous charge uniforme~:

\begin{equation}
w_{inst} = \frac{5 \, q \, L^4}{384 \, E_{0,mean} \, I}
\label{eq:winst}
\end{equation}

\begin{lstlisting}[style=python, caption={\fichier{ec5/els.py} --- flèche instantanée}]
def calculer_w_inst(q_kNm, L_m, E_mean_MPa, I_cm4) -> float:
    """Flèche instantanée w_inst = 5qL⁴/(384EI) en mm."""
    # Conversions d'unités :
    q_Nmm = q_kNm * 1.0     # kN/m = N/mm (numériquement identique !)
    L_mm  = L_m * 1000.0    # m -> mm
    I_mm4 = I_cm4 * 1e4     # cm⁴ -> mm⁴

    return 5.0 * q_Nmm * L_mm**4 / (384.0 * E_mean_MPa * I_mm4)
\end{lstlisting}

\begin{maitre}
Pourquoi \py{q\_Nmm = q\_kNm * 1.0}~? La conversion est triviale~?
\end{maitre}

\begin{apprenti}
Oui~: 1 kN/m $=$ 1000 N / 1000 mm $=$ 1 N/mm, le facteur est 1.
\end{apprenti}

\begin{maitre}
Exact. La multiplication par 1{,}0 est gard\'ee pour la
\textbf{lisibilit\'e}~: le commentaire montre qu'il y a bien une conversion
d'unit\'e, m\^eme si le facteur num\'erique est 1. C'est le Principe~VI
du projet~: toujours rendre les unit\'es explicites.
\end{maitre}

\subsection{Fl\`eche finale~: le fluage}

Le bois flue sous charge permanente.
La fl\`eche finale tient compte de ce fluage via le facteur $k_{def}$~:

\begin{equation}
w_{fin} = w_{inst} \cdot (1 + k_{def})
\label{eq:wfin}
\end{equation}

\begin{lstlisting}[style=python, caption={\fichier{ec5/els.py} --- flèche finale avec fluage}]
def calculer_w_fin(w_inst, k_def) -> float:
    """Flèche finale w_fin = w_inst × (1 + k_def)."""
    return w_inst * (1.0 + k_def)

# Valeurs typiques de k_def (EC5 Tableau 3.2) :
# Bois massif, classe de service 1 : k_def = 0.6
# Bois massif, classe de service 2 : k_def = 0.8
# Bois lamellé-collé, CS1           : k_def = 0.6
\end{lstlisting}

\subsection{Limites de fl\`eche et taux ELS}

Les limites admissibles sont lues depuis \fichier{limites\_fleche\_ec5.csv}~:

\begin{center}
\begin{tabular}{lrrr}
\toprule
\textbf{Usage} & \boldmath{$w_{inst} \leq$} & \boldmath{$w_{fin} \leq$} & \boldmath{$w_2 \leq$}\\
\midrule
Toiture inaccessible & $L/150$ & $L/100$ & ---\\
Plancher habitation   & $L/300$ & $L/200$ & $L/300$\\
Plancher bureaux      & $L/300$ & $L/250$ & $L/350$\\
\bottomrule
\end{tabular}
\end{center}

Le taux ELS est calcul\'e de fa\c{c}on analogue au taux ELU~:

\begin{equation}
\text{taux}_{ELS,inst} = \frac{w_{inst}}{w_{inst,lim}}
\leq 1{,}0
\end{equation}

\begin{lstlisting}[style=python, caption={Extrait de \fichier{ec5/els.py} --- boucle de vérification ELS}]
for combi in combinaisons:
    if combi.type_combinaison not in ("ELS_CAR", "ELS_FREQ", "ELS_QPERM"):
        continue  # on ne calcule les flèches QUE pour les combi ELS

    # Charges via polymorphisme (comme pour l'ELU)
    charges = type_poutre.charges_lineaires(config, materiau, longueurs_m, combi)

    for i, L_m in enumerate(longueurs_m):
        q_d = float(charges["q_d_kNm"][i])

        # Limites depuis le CSV normatif
        limites = get_limites_fleche(config.usage, L_m)
        limite_inst_mm = limites["limite_inst_mm"]
        limite_fin_mm  = limites["limite_fin_mm"]

        # Calcul des flèches
        w_inst_mm  = calculer_w_inst(q_d, L_m, materiau.E_0_mean_MPa, materiau.I_cm4)
        w_creep_mm = k_def * w_inst_mm      # composante de fluage
        w_fin_mm   = calculer_w_fin(w_inst_mm, k_def)

        # Taux de sollicitation ELS
        taux_inst = w_inst_mm / limite_inst_mm
        taux_fin  = w_fin_mm  / limite_fin_mm
\end{lstlisting}